\section{Problem Formulation}
\subsection{Known Directed Road Routing}
Let $G=(V,E,w)$ be a finite directed road graph with finite nonnegative weights $w:E\to\mathbb{R}_{\geq0}$. For origin $s\in V$ and destination $t\in V$, the routing objective is
\begingroup\predisplaypenalty=0
\begin{equation}
P^*\in\underset{P:s\leadsto t}{\operatorname{argmin}}\;
C(P),\quad C(P)=\sum_{(u,v)\in P}w(u,v).
\end{equation}
\endgroup
Here $C(P)$ sums exported road lengths in meters. Directed edges enforce one-way topology; traffic, time-dependent costs, and turn restrictions absent from preprocessing are excluded. If no path exists, search reports failure. Node endpoints isolate search from coordinate snapping in the principal experiment. Interactive requests instead use projected road positions and temporary partial edges (Section~4).

\subsection{Partially Observed Robot Navigation}
Let $W\subset\mathbb{R}^2$ be a bounded workspace and $O\subset W$ the static polygonal obstacles. At controller decision $k$, define robot position $x_k$, goal $q$, accumulated observed free space $F_k$, and visible obstacle fragments $B_k$. The simulator uses $O$ for sensing and motion validation; the policy receives observed geometry and traversal history. Its unknown region is $U_k=W\setminus(F_k\cup B_k)$. Range $R>0$ and occlusion limit observations: membership in the sensing disk alone establishes neither visibility nor traversability.

The controller selects target $z_k$ from observed information and history, then constructs $G_k=(V_k,E_k)$ from scan centers and request endpoints. Only wholly verified free segments become links. Edge weights $\|u-v\|_2$ and heuristic $\|u-z_k\|_2$ let \astar{} minimize graph-path length. Execution checks collision and observed-free coverage. For consecutive executed positions $y_0,\ldots,y_M$, the complete journey has length
\begin{equation}
L_{\mathrm{online}}=\sum_{j=0}^{M-1}\|y_{j+1}-y_j\|_2.
\end{equation}
This quantity includes both exploration and return motion. The index $j$ counts executed segments, whereas $k$ counts controller decisions; neither counts rendered frames. Consequently, $L_{\mathrm{online}}$ need not equal one local plan's cost or the shortest distance available with full-map knowledge.

The continuous simulator uses an idealized point center and reversible motion along straight segments. Heading changes freely at vertices; acceleration, nonholonomic constraints, and disk clearance are absent. The frozen local comparison permits boundary tangency in its primary convention and separately reports strict no-contact validity. Neither convention validates the source's nominal radius-0.5 parameter.

\subsection{Recovery and Evaluation Objectives}
A branch records an observed anchor and its subsequent executed trajectory. Exhaustion of usable local candidates triggers return to an observed ancestor with remaining alternatives. Failure to reach one target instead defers it for reconsideration after knowledge changes; it does not certify branch exhaustion. Exploration resumes after return. Section~5 derives the conditional executed-length bound.

Road evaluation measures minimum graph cost, valid expansions, and search-call time. Local robot evaluation measures planned Euclidean length, geometric turns, endpoint agreement, observed-space coverage, obstacle validity, and computation components. Supporting online studies measure actual distance, goal status, frontier processing, and executed returns. A return to an observed ancestor is \emph{certified} only when the recorded entry trajectory is reversible and the executed return satisfies the prescribed endpoint and length constraints. Reaching a goal, ending with no eligible frontier, and exhausting a budget are different outcomes; none automatically proves a complete map or global unknown-world optimality.
